% Part:sets-functions-relations
% Chapter: size-of-sets
% Section: enumerations-alt

\documentclass[../../../include/open-logic-section]{subfiles}

\begin{document}

\olfileid{sfr}{siz}{enm-alt}

\olsection{Opsommingen en \usetoken{S}{enumerable} verzamelingen}

\begin{editorial}
  Deze paragraaf definieert opsommingen als bijecties met
  (beginstukken van) $\Nat$, zoals gebruikelijk in de
  verzamelingenleer. Daardoor wijkt zij enigszins af van de definities
  in \olref[enm]{sec} en herhaalt zij alle voorbeelden daar. Zij is
  ook wat beknopter dan die paragraaf.
\end{editorial}

We kunnen een eindige verzameling eenvoudig opgeven door haar
!!{element}s op te sommen. Dat doen we wanneer we een verzameling als
volgt definiëren:
\[
  A = \{a_1, a_2, \ldots, a_n\}.
\]
Als de !!{element}s $a_1$, \dots, $a_n$ allemaal verschillend zijn,
levert dit !!a{bijection} op tussen $A$ en de eerste $n$ natuurlijke
getallen $0$, \dots, $n-1$. Omgekeerd heeft iedere eindige
verzameling slechts eindig veel !!{element}s en kan zij dus in zo'n
correspondentie worden gebracht. Met andere woorden: als $A$ eindig
is, bestaat er !!a{bijection} tussen $A$ en $\{0, \dots, n-1\}$,
waarbij $n$ het aantal !!{element}s van~$A$ is.

Als we bepaalde soorten oneindige verzamelingen toestaan, kunnen ook
sommige oneindige verzamelingen worden opgesomd. Dit kunnen we precies
formuleren door te zeggen dat een oneindige verzameling wordt
opgesomd door !!a{bijection} tussen die verzameling en heel~$\Nat$.

\begin{defn}[Opsomming, verzamelingstheoretisch]
Een \emph{opsomming} van een verzameling $A$ is !!a{bijection} met
beeld $A$ en met als domein een beginstuk van de natuurlijke getallen
$\{0, 1, \ldots, n\}$ {of} de gehele verzameling natuurlijke
getallen~$\Nat$.
\end{defn}

\begin{explain}
Het woord \emph{opsomming} past intuïtief bij deze definitie. Zeggen
dat we een verzameling $A$ hebben opgesomd, betekent immers dat er
!!a{bijection} $f$ bestaat waarmee we de elementen van de verzameling
$A$ één voor één kunnen tellen. Het $0$-de element is $f(0)$, het
eerste is $f(1)$, \ldots het $n$-de is $f(n)$\ldots.\footnote{We
tellen inderdaad vanaf $0$. We zouden uiteraard ook bij~$1$ kunnen
beginnen. Dat zou nauwelijks verschil maken: we zouden alleen~$\Nat$
door~$\PosInt$ moeten vervangen.} De reden voor deze definitie wordt
nog duidelijker met de volgende aanvulling:
\end{explain}

\begin{defn}
  \ollabel{defn:enumerable}
  Een verzameling~$A$ is !!{enumerable} dan en slechts dan als $A =
  \emptyset$ of er een opsomming van~$A$ bestaat. We noemen $A$
  !!{nonenumerable} dan en slechts dan als $A$ niet !!{enumerable} is.
\end{defn}

\begin{explain}
Een verzameling is dus !!{enumerable} dan en slechts dan als zij leeg
is of als haar !!{element}s met een opsomming één voor één kunnen
worden geteld.
\end{explain}

\begin{ex}
De identiteitsfunctie $\Id{\Nat} \colon \Nat \to \Nat$, gegeven door
$\Id{\Nat}(n) = n$, somt de natuurlijke getallen op. De functie die de
\emph{positieve} natuurlijke getallen $\Nat^+ = \Nat \setminus
\{0\}$ opsomt, is de functie $g(n) = n + 1$, dus de opvolgerfunctie.
\end{ex}

\begin{prob}
Bewijs dat een verzameling $A$ !!{enumerable} is dan en slechts dan
als $A = \emptyset$ of er !!a{surjection} $f\colon \Nat \to A$
bestaat. Bewijs dat $A$ !!{enumerable} is dan en slechts dan als er
!!a{injection} $g\colon A \to \Nat$ bestaat.
\end{prob}

\begin{ex}
De functies $f\colon \Nat \to \Nat$ en $g \colon \Nat \to \Nat$,
gegeven door
\begin{align*}
  f(n) & = 2n \text{ en}\\
  g(n) & = 2n+1
\end{align*}
sommen respectievelijk de even en de oneven natuurlijke getallen op.
Geen van beide is echter !!{surjective}; geen van beide is dus een
opsomming van $\Nat$.
\end{ex}

\begin{prob}
Definieer een opsomming van de kwadraatgetallen $1$, $4$, $9$, $16$,
\dots
\end{prob}

\begin{ex}
Laat $\lceil x \rceil$ de \emph{plafondfunctie} zijn, die $x$ naar
boven afrondt op het dichtstbijzijnde gehele getal. Dan somt de functie
$f \colon \Nat \to \Int$, gegeven door
\[
  f(n) = (-1)^{n} \left\lceil\tfrac{n}{2}\right\rceil
\]
de gehele getallen~$\Int$ als volgt op:
\[
\begin{array}{c c c c c c c c}
f(0) & f(1) & f(2) & f(3) & f(4) & f(5) & f(6) & \dots \\ \\
\big\lceil \tfrac{0}{2} \big\rceil & -\big\lceil \tfrac{1}{2}\big\rceil &  \big\lceil \tfrac{2}{2} \big\rceil & -\big\lceil \tfrac{3}{2} \big\rceil & \big\lceil \tfrac{4}{2} \big\rceil  & -\big\lceil \tfrac{5}{2}\big\rceil & \big\lceil \tfrac{6}{2} \big\rceil & \dots \\ \\
0 & -1 & 1 & -2 & 2 & -3 & 3& \dots
\end{array}
\]
Merk op hoe $f$ de waarden van $\Int$ geeft door heen en weer te
‘springen’ tussen positieve en negatieve gehele getallen. We kunnen
$f$ ook als volgt stuksgewijs definiëren:
\[
f(n) = \begin{cases}
  \frac{n}{2} & \text{als $n$ even is}\\
  -\frac{n+1}{2} & \text{als $n$ oneven is}
  \end{cases}
\]
\end{ex}

\begin{prob}
Bewijs dat als $A$ en $B$ !!{enumerable} zijn, ook $A \cup B$ dat is.
\end{prob}

\begin{prob}
Bewijs met inductie naar $n$ dat als $A_1$, $A_2$, \dots, $A_n$
allemaal !!{enumerable} zijn, ook $A_1 \cup \dots \cup A_n$ dat is.
\end{prob}
\end{document}